\documentclass[11pt,a4paper]{article}
\usepackage[T1]{fontenc}
\usepackage[utf8]{inputenc}
\usepackage[ngerman]{babel}
\usepackage[margin=2.5cm]{geometry}
\usepackage[hidelinks]{hyperref}
\usepackage{array}
\usepackage{forest}
\useforestlibrary{edges}
\setlength{\parindent}{0pt}
\setlength{\emergencystretch}{2em}
\renewcommand{\arraystretch}{1.25}
% Zeilenumbruch und Leerzeichen in Textknoten
\newcommand{\RET}{\ensuremath{\hookleftarrow}\allowbreak}
\newcommand{\VS}{\textvisiblespace}
\newcommand{\IV}{\VS\VS\VS\VS\allowbreak}
\forestset{
  el/.style={draw=blue!60!black, fill=blue!6, rounded corners=1.5pt},
  tx/.style={draw=black!50, dashed, fill=black!4},
  ot/.style={draw=black!60, dotted},
  at/.style={label={[font=\ttfamily\scriptsize, text=black!60, anchor=west, xshift=2pt]east:{#1}}},
}
\begin{document}
\shorthandoff{"}
\begin{center}
  {\Large\textbf{Webtechnologien -- Übungsblatt 2}}
\end{center}
\section*{Aufgabe 3: DOM Operations}
\textbf{Aufgabe:} Den DOM von \texttt{ex3.html} skizzieren, den
vorgegebenen JavaScript-Code einbinden und jede Konsolenausgabe erklären.
Zusätzlich die Wirkung von \texttt{removeAttribute("target")} auf den
dritten Link beschreiben.

\medskip
\textbf{3.1 DOM-Skizze:} Originale \texttt{ex3.html} ohne das eingefügte
\texttt{script}-Element, Textknoten im Browser nachgeprüft.

\begin{center}
\begin{forest}
  for tree={folder, grow'=0, font=\ttfamily\small, inner sep=1.5pt,
    l sep=7pt, s sep=4.5pt, edge={black!55}}
  [\#document, ot
    [<!DOCTYPE html>, ot]
    [html, el, at={lang="en"}
      [head, el
        [\#text "\RET\IV", tx]
        [meta, el, at={charset="UTF-8"}]
        [\#text "\RET\IV", tx]
        [title, el
          [\#text "DOM Operations", tx]]]
      [\#text "\RET", tx]
      [body, el
        [\#text "\RET", tx]
        [a, el, at={href="https://en.wikipedia.org" target="\_new"}
          [img, el, at={src="\dots" alt="British Flag"}]]
        [\#text "\RET BeforeParent\RET", tx]
        [a, el, at={href="https://de.wikipedia.org" target="\_new"\quad{\rmfamily\itshape\normalsize$\leftarrow$ link in 3.2}}
          [p, el
            [\#text "Text", tx]]
          [\#text "\RET\IV Before\VS", tx]
          [img, el, at={src="\dots" alt="Austrian Flag"}]
          [\#text "\RET\IV After\RET", tx]]
        [\#text "\RET AfterParent\RET", tx]
        [a, el, at={href="https://ja.wikipedia.org" target="\_new"\quad{\rmfamily\itshape\normalsize$\leftarrow$ link in 3.3}}
          [\#text "\RET\IV", tx]
          [img, el, at={src="\dots" alt="Japanese Flag"}]
          [\#text "\RET", tx]]
        [\#text "\RET\RET", tx]]]]
\end{forest}
\end{center}

{\small Durchgezogener Rahmen: Elementknoten; gestrichelt: Textknoten
(\RET{} = Zeilenumbruch, \VS{} = Leerzeichen); gepunktet: Dokument- und
Doctype-Knoten. Grau rechts: Attribute. Sie gehören zum Element, sind aber
keine Kindknoten. Der Zeilenumbruch zwischen \texttt{</head>} und
\texttt{<body>} landet als Textknoten im \texttt{html}-Element, die
Zeilenumbrüche nach \texttt{</body>} im \texttt{body}.}

\medskip
\textbf{3.2 Konsolenausgaben:} \texttt{dom.js} ist im \texttt{head} mit
\texttt{<script src="dom.js" defer>} eingebunden. Durch \texttt{defer}
läuft das Skript erst, wenn der DOM vollständig aufgebaut ist.
\texttt{link = document.getElementsByTagName("a")[1]} ist der
\emph{zweite} Link (\texttt{de.wikipedia.org}), da ab 0 gezählt wird.
Zeichenketten sind hier mit \RET{} und \VS{} dargestellt.

{\small
\begin{center}
\begin{tabular}{|>{\raggedright\arraybackslash}p{0.45cm}|>{\raggedright\arraybackslash}p{4.45cm}|>{\raggedright\arraybackslash}p{4.0cm}|>{\raggedright\arraybackslash}p{5.25cm}|}
\hline
 & \textbf{Ausdruck} & \textbf{Ausgabe} & \textbf{Begründung} \\ \hline
1 & \nolinkurl{link.attributes["href"].nodeValue} & \nolinkurl{https://de.wikipedia.org} & Wert des Attributknotens \texttt{href}, genau wie im HTML geschrieben. \\ \hline
2 & \nolinkurl{link.getAttributeNode("href").nodeValue} & \nolinkurl{https://de.wikipedia.org} & Derselbe Attributknoten, nur über eine Methode geholt. \\ \hline
3 & \nolinkurl{link.href} & \nolinkurl{https://de.wikipedia.org/} & DOM-Property: aufgelöste, normalisierte URL, daher mit \texttt{/} am Ende. \\ \hline
4 & \nolinkurl{link.firstChild} & Element \texttt{<p>Text</p>} & \texttt{<p>} steht direkt nach \texttt{<a ...>}, davor gibt es keinen Textknoten. \\ \hline
5 & \nolinkurl{link.attributes[0].nodeType} & \texttt{2} & \texttt{attributes[0]} ist \texttt{href}, ein Attributknoten (\texttt{ATTRIBUTE\_NODE}). \\ \hline
6 & \nolinkurl{link.firstChild.nodeType} & \texttt{1} & Das \texttt{p} ist ein Elementknoten (\texttt{ELEMENT\_NODE}). \\ \hline
7 & \nolinkurl{link.parentNode.nodeName} & \texttt{BODY} & Der Link liegt direkt im \texttt{body}; \texttt{nodeName} von HTML-Elementen ist großgeschrieben. \\ \hline
8 & \nolinkurl{link.lastChild.parentNode.nodeName} & \texttt{A} & \texttt{lastChild} ist der Textknoten \texttt{"\RET\IV After\RET"}, sein Elternknoten wieder der Link. \\ \hline
9 & \nolinkurl{link.attributes.length} & \texttt{2} & Der Link hat zwei Attribute: \texttt{href} und \texttt{target}. \\ \hline
10 & \nolinkurl{link.innerHTML} & \texttt{<p>Text</p>\RET\IV Before\VS <img src="\dots" alt="Austrian Flag">\RET\IV After\RET} & HTML-Inhalt innerhalb des \texttt{a} als String; Whitespace bleibt erhalten. \\ \hline
11 & \nolinkurl{link.outerHTML} & \texttt{<a href="https://\allowbreak de.wikipedia.org" target="\_new">}, dann wie 10, dann \texttt{</a>} & Wie \texttt{innerHTML}, aber inklusive des \texttt{a}-Tags selbst. \\ \hline
12 & \nolinkurl{link.innerText} & \texttt{Text\RET\RET Before\VS\VS After\VS} & Sichtbarer Text wie gerendert: Whitespace zusammengefasst, nach dem Blockelement \texttt{p} zwei Umbrüche. Links und rechts vom Bild bleibt je ein Leerzeichen. Am Ende ein Leerzeichen, weil danach \glqq AfterParent\grqq{} in derselben Zeile folgt. \\ \hline
13 & \nolinkurl{link.outerText} & wie 12 & Beim Lesen gleich wie \texttt{innerText}. Beim Setzen ersetzt \texttt{outerText} das ganze Element, \texttt{innerText} nur den Inhalt. \\ \hline
14 & \nolinkurl{link.textContent} & \texttt{Text\RET\IV Before\VS\RET\IV After\RET} & Alle Textknoten roh hintereinander, ohne Layout. Der \texttt{alt}-Text ist ein Attribut und fehlt daher. \\ \hline
15 & \nolinkurl{link} & Element \texttt{<a href="\dots" target="\_new">\dots</a>} & Das Elementobjekt selbst, in der Konsole aufklappbar. \\ \hline
16 & \nolinkurl{link.nodeValue} & \texttt{null} & Elementknoten haben keinen \texttt{nodeValue}, nur Text-, Attribut- und Kommentarknoten. \\ \hline
\end{tabular}
\end{center}}

{\small Bei 10, 11 und 15 ist die Bild-URL mit \glqq\dots\grqq{} gekürzt,
in der Konsole steht die vollständige URL. Knotentypen: 1 Element,
2 Attribut, 3 Text, 8 Kommentar, 9 Dokument.}

\medskip
\textbf{3.3 \texttt{removeAttribute}:}
\texttt{document.getElementsByTagName("a").item(2)} ist der \emph{dritte}
Link (\texttt{ja.wikipedia.org}); \texttt{item(2)} ist gleichwertig zu
\texttt{[2]}. \texttt{link.removeAttribute("target")} entfernt nur bei
diesem Link das Attribut \texttt{target="\_new"}. \texttt{\_new} ist kein
reserviertes Ziel wie \texttt{\_blank}, sondern ein Fenstername: Der erste
Klick öffnet einen neuen Tab mit diesem Namen, weitere Links mit
\texttt{target="\_new"} laden in denselben Tab. Ohne \texttt{target}
öffnet der japanische Link im aktuellen Tab, da kein \texttt{<base target>}
gesetzt ist. Linkziel und Bild bleiben gleich, die anderen beiden Links
behalten \texttt{target="\_new"}. Sichtbar ändert sich auf der Seite nichts.

\medskip
\textbf{Hinweis:} \texttt{var link} wird in 3.3 ein zweites Mal
deklariert. Mit \texttt{var} ist das erlaubt, mit \texttt{let} wäre es
ein \texttt{SyntaxError}.

\medskip
\textbf{Test:} In Chrome, \texttt{ex3.html} direkt geöffnet. Die 16
Ausgaben in der Konsole entsprechen der Tabelle. Die \texttt{target}-Werte
der drei Links sind danach \texttt{"\_new"}, \texttt{"\_new"} und
\texttt{null}.

\medskip
{\small \textbf{Hilfsmittel:} Claude Code (Claude Opus 5.5) für Erstellung, Prüfung und Dokumentation.}
\end{document}
